\scanpage{33}

\uline{Fermés de Lefschetz-Grauert}

1) Définitions, remarques

\tikz[remember picture,overlay]{\coordinate (lgda);}\uline{Déf} Soit $X$ un préschéma loc noeth, $Y$ une partie fermée, $n$ un entier. On dit que $Y$ a la propriété $(LG)_n$ dans $X$ \deleted{\unsure{jusqu'en|presque} \unsure{dim|si et|dans} \unsure{$n$|un}}, si pour tt $F$ loc. libre sur $X$,

\[
H^i(X,F)\mathbin{\unsure{-|\to}}H^i(X_{/Y},F_{/Y})\deleted{\unsure{\mathrm{est}|\mathrm{isom}}}
\tag*{\makebox[0pt][r]{\makebox[\linewidth][l]{$(\ast)$}}}
\]

est bijectif pour $i\leq n$, \inserted{et injectif pour $i=n+1$.}\margin{inutile}

On dit que $Y$ a la propriété $LG$ \inserted{faible dans $X$} si \deleted{elle a la propriété $(LG)_0$,} \inserted{pour} tout $F$ loc. libre sur $X$,

\[
H^0(X,F)\overset{\sim}{\longrightarrow}H^0(X_{/Y},F_{/Y})
\]

\tikz[remember picture,overlay]{\coordinate (lgdb); \draw ([xshift=-6pt,yshift=-3pt]lgda) -- ([xshift=-6pt]lgda |- lgdb); }[$i.e.$ $F\leadsto F_{/Y}$ est ple. fid. des Mod. loc. libres sur $X$ en ceux sur $X_{/Y}$]

On dit que $Y$ a la propriété $LG$ \deleted{\unsure{stricte|de stricte}} \inserted{dans $X$} si tt \unsure{pour|$V$} voisinage ouvert $U$ de $Y$, $Y$ \deleted{ait} \inserted{a la} propriété $LG$ \inserted{faible dans $U$}. \uline{\unsure{Enfin}} $Y$ a la propriété $LG$, \deleted{\unsure{stricte}} \inserted{\unsure{définitive|effective}} si elle a la propriété $LG$, \inserted{et si de plus} \deleted{pour}
\margin{$\Rightarrow$ Si $z\in X$ est tel \deleted{\unsure{fermé et ne ren-|qu'il se ren-}} \deleted{\unsure{contre pas $Y$|contre pas}} que $\bar z\cap Y=\varnothing$, alors $z\notin\operatorname{Ass}\underline{\mathcal O}_X$.}
tt $\mathcal F$ loc libre sur $X_{/Y}$ \deleted{est} est isom : un faisceau de la forme $F_{/Y}$, où $F$ loc libre sur un voisinage ouvert $U$ de $Y$.
% slanted
\margin{Question : $(\unsure{\mathrm{Lef}|\mathrm{LGf}|\mathrm{LG}})$ \deleted{\unsure{\#}} implique-t-elle \deleted{\unsure{formellement}} $\forall x\in X$ tq $\bar x\cap Y=\varnothing$, on a $\operatorname{prof}\mathcal O_{X,x}\geq 2$ ???}

Donc la propriété \textbf{LG} stricte signifie que les foncteurs $\mathcal L(U)\mathbin{-}\mathcal L(\unsure{W|V})$ $[\deleted{\unsure{\forall}}U\supset V\supset Y]$ sont pleinement fidèles, et

\[
\underset{U\supset Y}{\underset{\longrightarrow}{\tikz[remember picture,baseline=(lglim.base)]{\node[inner sep=0pt](lglim){$\operatorname{Lim}$};}}}\;
\underbrace{\mathcal L(U)}_{\substack{\text{Catégorie des}\\\text{Modules loc}\\\text{libres sur }U}}
\longrightarrow\mathcal L(X_{/Y})
\]
\tikz[remember picture,baseline=(lgcaption.base)]{\node[align=center,inner sep=1pt](lgcaption){``Pseudolim'' de\\catégories};}
\tikz[remember picture,overlay]{\draw[<->] (lgcaption.north east) -- (lglim.south west);}

est une équivalence de catégories.
% slanted
\margin{\uline{Exemples tout de suite}}

% archive sheet number: 32
